\section{Warmup: cái vá cho chuyện vừa đo được}
\label{sec:ch08-warmup}

\index{warmup!công thức của mục 5.3}%
Mục 5.3 của bài dài đúng bốn câu và một phương trình. Adam với
\(\beta_1 = 0.9\), \(\beta_2 = 0.98\), \(\epsilon = 10^{-9}\), rồi:

\begin{equation}
  \mathit{lrate} = \dmodel^{-0.5} \cdot
  \min\!\left( s^{-0.5},\; s \cdot w^{-1.5} \right),
  \label{eq:ch08-warmup}
\end{equation}

với \(s\) là số thứ tự bước và \(w\) là \code{warmup\_steps}, mà bài đặt bằng
4000. Bài giải thích bằng một câu: \enquote{This corresponds to increasing the
learning rate linearly for the first \(warmup\_steps\) training steps, and
decreasing it thereafter proportionally to the inverse square root of the step
number.}

Đọc lướt thì đây là ba hằng số gõ tay. Viết ra thì nó có ba tính chất, và cả ba
đều dẫn được ra từ chính \eqref{eq:ch08-warmup} chứ không cần thí nghiệm nào.

\index{warmup!ba tính chất dẫn ra được}%
\textbf{Hai nhánh gặp nhau đúng ở đỉnh.} Ở \(s = w\), nhánh trái cho
\(w^{-0.5}\) và nhánh phải cho \(w \cdot w^{-1.5} = w^{-0.5}\). Bằng nhau, nên
\(\min\) không tạo ra một bước nhảy: lịch này là một đoạn thẳng nối liền vào một
đường cong, và \(w\) là hoành độ đỉnh chứ không phải một chỗ chuyển chế độ.

\textbf{Đỉnh cao bao nhiêu thì đọc ra ngay.} Thay vào:
\(\dmodel^{-0.5} w^{-0.5} = (\dmodel \, w)^{-1/2}\). Nên learning rate lớn nhất
mà lịch này từng đạt tới do \emph{tích} của bề rộng và số bước warmup quyết
định, và hai tham số ấy đổi nhau được.

\textbf{Mô hình rộng hơn nhận learning rate nhỏ hơn.} Thừa số
\(\dmodel^{-0.5}\) nằm ngoài cả \(\min\), nên đi từ \(\dmodel = 64\) lên 512 là
chia learning rate cho \(\sqrt{8}\). Đó là lý do một con số warmup chép từ bài
này sang mô hình khác bề rộng thì không còn nghĩa gì, và
\code{tests/test\_ch08.py} khẳng định đúng tỷ số ấy.

Trong repo, phần tính của \eqref{eq:ch08-warmup} là bảy dòng, chép ở đây không
kèm docstring và không kèm dòng kiểm \code{warmup\_steps} phải dương:

\begin{minted}{python}
def warmup_lambda(d_model: int, warmup_steps: int):
    def scale(step: int) -> float:
        step_num = step + 1
        return d_model ** -0.5 * min(
            step_num ** -0.5, step_num * warmup_steps ** -1.5
        )

    return scale
\end{minted}

Cái \code{step + 1} không phải trang trí. \code{LambdaLR} đếm từ 0 và
\eqref{eq:ch08-warmup} chia cho \(s\), nên để nguyên thì bước đầu tiên chia cho
không.

\subsection*{Vì sao nó là một cái vá, chứ không phải một mẹo}

\index{warmup!nối với gradient lúc khởi tạo}%
Đặt \eqref{eq:ch08-warmup} cạnh bảng ở mục~\ref{sec:ch08-thu-tu-ln} thì hai thứ
khớp vào nhau.

Ở lúc khởi tạo, gradient của post-LN tại tầng cuối lớn gấp khoảng năm lần của
pre-LN, và không nhỏ đi khi chồng tầng dày lên. Một learning rate cố định nhân
với một gradient lớn ở bước đầu tiên là một bước cập nhật lớn ở đúng lúc mô hình
chưa biết gì. \eqref{eq:ch08-warmup} làm learning rate nhỏ gần bằng không ở
những bước ấy rồi mới cho lớn dần.

Nói cách khác, warmup không phải một cách để huấn luyện nhanh hơn. Nó là cách để
sống sót qua vài trăm bước đầu của một kiến trúc mà thứ tự layer norm đã làm cho
gradient ban đầu lớn. Ai đổi thứ tự thì bớt cần nó, và đó đúng là câu Xiong và
cộng sự rút ra.

\index{warmup!ba bài không nói cùng một câu}%
Chỗ này có một chuyện về nguồn đáng nói riêng, vì cách kể lại thường gộp ba bài
làm một. \enquote{Các bài pre-norm} không cùng nói rằng bỏ được warmup.

\begin{itemize}
  \item Wang và cộng sự, ACL 2019, \emph{không} nói thế~\autocite{wang2019deep}.
    Bài ấy dùng warmup và còn nâng lên 16000 bước cho mô hình sâu; luận điểm của
    nó là pre-norm làm encoder hai mươi tới ba mươi tầng huấn luyện được, tức
    chuyện độ sâu.
  \item Nguyen và Salazar, IWSLT 2019, là chỗ nói sớm
    nhất~\autocite{nguyen2019tears}. Tóm tắt viết \enquote{pre-norm residual
    connections (PRENORM) and smaller initializations enable warmup-free,
    validation-based training with large learning rates}.
  \item Xiong và cộng sự 2020 là chỗ chứng minh, và tự gọi Nguyen và Salazar là
    công trình song song độc lập.
\end{itemize}

Nguyen và Salazar kèm một điều kiện mà chỗ trích lại hay bỏ, và nó nằm ngay
trang đầu: \enquote{we find it degrades base Transformer performance on WMT'14
English-German}. Nên phát biểu đủ là pre-norm bỏ được warmup, và trên một tập dữ
liệu lớn nó làm mô hình base kém đi.

\begin{bridge}{Popel và Bojar 2018: đặt warmup thấp quá thì không phải chậm, mà là hỏng}
\index{warmup!Popel 2018}%
Bài này là một tập ghi chép vận hành về Tensor2Tensor, không phải một bài lý
thuyết~\autocite{popel2018tips}. Mục 4.6 quét \code{warmup\_steps} quanh mặc
định 16000 ở learning rate cố định, và kết quả đáng nhớ nhất là một chỗ hỏng
chứ không phải một đường cong: đặt 12000 thì huấn luyện \emph{phân kỳ}
(\enquote{Setting warmup steps too low (12k) results in diverged training}),
còn 48000 chỉ hội tụ chậm hơn lúc đầu rồi bắt kịp.

Kết luận của bài thì rộng rãi: \enquote{there is a relatively large range of
learning rate and warmup steps values that achieve the optimal results}. Hai vế
ấy đi cùng nhau mới đủ. Có một dải rộng chấp nhận được, và ngay dưới dải ấy
không phải là kém đi một chút mà là không huấn luyện được.

Đây là bằng chứng độc lập cho chuyện mục này nói, ở quy mô mà sách này không
chạm tới được. Tôi đọc bài ở mục 4.5 và 4.6.
\end{bridge}

\subsection*{Đo ở quy mô của sách, và nó không trả lời được}

\index{warmup!bảng bốn ô ở quy mô sách}%
Bây giờ đến phần tôi phải cẩn thận. Câu hỏi tự nhiên là chạy bốn cấu hình, hai
thứ tự nhân hai lịch, trên corpus của sách rồi đọc ô nào cao nhất.

Chương~\ref{ch:transformer} đã cảnh báo một nửa chuyện: nó thử lịch của bài ở
658 bước và thấy thua recipe dùng chung, rồi tự nói đó là kết quả về 658 bước.
Nửa còn lại là chuyện \tn{khớp đúng cả câu}{exact match} trên 300 câu là một
thống kê nhảy, và
mục~\ref{sec:ch07-corpus} đã in một trường hợp nó đi xuống khi số epoch đi lên.
Nên bảng này chạy ba seed cho mỗi ô, và cột \code{min} với \code{max} là hai cột
đáng đọc trước.

\begin{measured}{hai thứ tự nhân hai lịch, ba seed mỗi hàng}
\begin{minted}[fontsize=\footnotesize]{text}
order    schedule  loss      exact     min      max      short   long
post-LN  flat      0.1449    0.5333    0.4700   0.5667   0.8964  0.1798
post-LN  warmup    0.1484    0.4967    0.3800   0.6500   0.8581  0.1447
pre-LN   flat      0.1039    0.5989    0.5067   0.6733   0.9640  0.2434
pre-LN   warmup    0.0855    0.6611    0.5667   0.7800   0.9685  0.3618
\end{minted}

2 tầng, 14 epoch, \(\dmodel = 24\), \(w = 200\), seed 0, 1, 2.

Đo bằng \code{experiments/ch08\_norm.py} tại tag \code{ch08}.
\end{measured}

\index{warmup!khoảng trong một ô rộng hơn khoảng giữa các ô}%
Đọc cột \code{min} và \code{max} trước, và bảng đổi nghĩa. Hàng
\code{post-LN warmup} đi từ 0.3800 tới 0.6500 chỉ vì đổi seed, rộng 0.27. Khoảng
cách giữa trung bình cao nhất và thấp nhất trong cả bốn hàng là 0.6611 trừ
0.4967, rộng 0.16.

\emph{Khoảng biến động bên trong một cấu hình rộng hơn khoảng cách giữa các cấu
hình.} Nên bảng này không tách được tác dụng của lịch warmup ở quy mô này, và
tôi nói thẳng ra thay vì chọn ô cao nhất rồi kể một câu chuyện quanh nó.

Tôi biết chính xác cái bảng ấy sẽ nói gì nếu chạy một seed, vì tôi đã chạy một
seed trước. Nó cho 0.4700, 0.4600, 0.5067 và 0.6367, và từ bốn số đó tôi đã suýt
viết rằng warmup không giúp gì cho post-LN mà giúp nhiều cho pre-LN. Ba seed cho
thấy đó là seed chứ không phải lịch. Chương~\ref{ch:mot-vector-cho-moi-tu} gặp
đúng bài học này ở một tầng khác, khi một \tn{phép cắt bỏ}{ablation} không có
đối chứng cùng số tham số đo ra một hiệu ứng không có thật.

\index{warmup!cái vượt được nhiễu}%
Cái \emph{có} vượt được nhiễu thì cũng đáng ghi, và nó không phải chuyện lịch.
Hai hàng pre-LN có mất mát thấp hơn hai hàng post-LN, 0.1039 và 0.0855 so với
0.1449 và 0.1484, và cột câu dài của chúng cao hơn hẳn, 0.2434 và 0.3618 so với
0.1798 và 0.1447. Chiều thì nhất quán qua cả bốn ô. Độ lớn thì bảng này không
nói được.

Một kiểm chứng cuối, và nó là lý do hàng đầu tiên tồn tại. Hàng
\code{post-LN flat} ở riêng seed 0 cho \code{0.1159 / 0.4700 / 0.8581 / 0.0921},
trùng từng chữ số với bảng ở mục~\ref{sec:ch07-corpus}. Chương này thêm hai tùy
chọn vào lớp Transformer của repo, và nếu chúng lỡ đụng vào đường mặc định thì
hàng ấy đã lệch.

\index{warmup!vì sao quy mô này không đủ}%
Còn vì sao quy mô này không trả lời được thì nằm ngay trong định lý. Chặn của
pre-LN nhỏ đi theo \(1/\sqrt{L}\), nên chuyện hai thứ tự khác nhau là chuyện của
\(L\) lớn. Ở đây \(L = 2\), và \(\sqrt{2}\) không phải một hệ số đáng đi tìm
trong một cột nhảy 0.27 vì seed. Bài chạy sáu tầng mỗi phía và 100000 bước cho
model base; Wang
và cộng sự phải lên hai mươi tầng mới thấy pre-norm đáng giá. Sách này chạy hai
tầng và 658 bước vì sách này hứa với bạn rằng mọi thí nghiệm chạy xong trên một
máy không có card đồ họa trong vài phút, và đây là chỗ lời hứa ấy phải trả giá.

Nên câu trung thực là: mục này dựng được nửa lý thuyết bằng phép đo tất định ở
lúc khởi tạo, và không dựng được nửa thực nghiệm. Bảng bốn ô ở trên là bằng
chứng cho việc nó không dựng được, chứ không phải một kết quả yếu.
